\section{Experiments}

\begin{figure}[t]
  \begin{tabular}{ccc}
    $\mathrm{P}[\dim(p)=1]$ &
    $\mathrm{P}[\dim(p)=2]$ &
    $\mathrm{P}[\dim(p)=3]$ \\
    \includegraphics[width=0.32\textwidth]{test0-n0-r5-dim1.png}&
    \includegraphics[width=0.32\textwidth]{test0-n0-r5-dim2.png}&
    \includegraphics[width=0.32\textwidth]{test0-n0-r5-dim3.png}\\
    \includegraphics[width=0.32\textwidth]{test0-n0_2-r5-dim1.png}&
    \includegraphics[width=0.32\textwidth]{test0-n0_2-r5-dim2.png}&
    \includegraphics[width=0.32\textwidth]{test0-n0_2-r5-dim3.png}\\
    %% \includegraphics[width=0.48\textwidth]{test0-n0-r4-dim.png}&
    %% \includegraphics[width=0.48\textwidth]{test0-n0-r7-dim.png}\\
    %% \includegraphics[width=0.48\textwidth]{test0-n0_2-r4-dim.png}&
    %% \includegraphics[width=0.48\textwidth]{test0-n0_2-r7-dim.png}\\
  \end{tabular}
  \caption{\label{fig-dimension-analysis}Illustration of the local
    $m$-measures to recognize the local dimension of a digital point
    in a toy 3D image, composed of line segments, pierced surfaces and
    a thick annulus with thin edges (top row: original image, bottom
    row: image perturbated by a Kanungo noise with parameter 0.2). The
    colormap ranges from white (0.0) to dark red (1.0). Each column
    gives the probability that a voxel has a given dimension $k$, for
    a radius $r=5$.}
\end{figure}

For a point $p$, let $W:=(W_0,\ldots,W_d)$ be its local measures for a
radius $r$. Then define
\[
d_{k,j} := \left\{
\begin{array}{l}
  (W_j - \bar{W}_j(k))/\bar{W}_j(k))~\text{when $0\Le j \Le k$}\\
  W_j/\bar{W}_j(j)~\text{when $k < j \Le d$}
\end{array}
\right.,
\qquad
q_k := \Pi_{j=0}^d e^{-j d_{k,j}^2}.
\]
The value $q_k$ is the greatest when the vector $W$ ressembles the
vector $\bar{W}(k)$. Then the probability that $p$ has dimension $k$
can be expressed as $\mathrm{P}[\dim(p)=k] := \mathrm{softmax}_k(q_0,
\ldots, q_d)$. \RefFigure{fig-dimension-analysis} illustrates these
values on a toy 3D image assembling digitized submanifolds of
different dimensions. Note how the probabilities are invariant to the
orientation of digitized manifolds, while staying robust to noise.
